\chapter{Variables and Memory}\label{ch:variables}

\begin{goals}
\begin{itemize}
  \item How the computer's memory holds things
  \item Declaring a variable with \code{:=} and changing its value with \code{=}
  \item The difference between a fixed value and a changeable one (\code{mut})
  \item The naming rules and how to choose a good name
  \item A few everyday patterns: the counter, adding up, swapping
\end{itemize}
\end{goals}

\section{Boxes with labels}

In the last chapter we wrote every number right where it was needed:
\code{println 2026 - 2009}. Real programs, however, have to \emph{remember}
things: the user's name, the amount of a purchase, the score in a game, the
number of times something has been done. The place where a program keeps
these things is the computer's \textbf{memory}.

Picture memory as a large shelf of small boxes. Each box can hold one thing:
a number, a piece of text, or other things we will meet later. So that we can
find the box again, we stick a label with a name on it. Such a box is called
a \textbf{variable}: a named place in memory that holds one value.

\section{Declaring a variable}

In Salam, to make a new box and put something in it, we use the sign
\code{:=}. Read it as ``colon equals'', or better still as ``define as'':

\program{The first variables}
\begin{salamcode}
func main:
    name := "Sara"
    age := 17
    println "Name:", name
    println "Age:", age
end
\end{salamcode}

\begin{salamout}
Name: Sara
Age: 17
\end{salamout}

The line \code{name := "Sara"} does three things: it makes a box, it sticks
the label \code{name} on it, and it puts the text \code{"Sara"} inside.
From then on, wherever we write \code{name} in the program, Salam looks in
that box and takes out its value. That is why \code{println "Name:", name}
printed ``Name: Sara'' and not ``Name: name''.

\begin{note}
Mind the difference between \code{"name"} and \code{name}. The first, with
quotation marks, is a piece of text that is printed just as it is. The
second, without quotation marks, is the name of a variable, and Salam puts
the value inside the box in its place.
\end{note}

Variables can be used in arithmetic, just like numbers:

\program{Arithmetic with variables}
\begin{salamcode}
func main:
    book price := 12
    count := 3
    println "Total:", book price * count, "dollars"
end
\end{salamcode}

\begin{salamout}
Total: 36 dollars
\end{salamout}

The advantage is clear: if the price of the book changes tomorrow, we change
one line, and every calculation that uses \code{book price} corrects itself.

\section{Naming variables}

In the program above one of the variables was called \code{book price}: two
words with a space between them. This is one of the pleasant features of
Salam: a variable name can be made of several words, the way we speak. The
naming rules are these:
\begin{itemize}
  \item A name may be made of letters, digits, the underscore \code{\_} and
        spaces.
  \item A name cannot start with a digit: \code{2people} is wrong, but
        \code{people2} is fine.
  \item A name cannot be one of the words of the language itself: you
        cannot make a variable called \code{println}, \code{end} or
        \code{if}. This rule applies to every word of a multi-word name
        too; because \code{in}, \code{to}, \code{by} and \code{each} are
        words of the language, names such as \code{items in cart} or
        \code{total to pay} are rejected. Appendix~\ref{app:keywords} lists
        these words.
  \item Upper and lower case are different; \code{x} and \code{X} are two
        separate variables.
\end{itemize}

More important than the rules are \emph{good habits}. The name of a variable
should say what is inside it. Compare:

\begin{salamcode}
func main:
    a := 12
    b := 3
    println a * b
end
\end{salamcode}

with the previous program. Both work out the same number, but the previous
program can still be understood a week later and this one cannot. For
temporary, unimportant names (say, in a two-line calculation) short names do
no harm, but for anything that stays in the program, choose a full,
descriptive name.

\begin{tip}
Salam lets you write names as single words joined with underscores as well,
\code{book\_price} for instance, and you will see that style in some of the
sample programs of the online editor. In this book we write multi-word names
with spaces, so that the programs read as plainly as possible.
\end{tip}

\section{Changing a value: fixed and changeable}

Now a point that sets Salam apart from many languages. A box you make with
\code{:=} is \textbf{sealed} by default: the value you first put in it stays
there to the end. Trying to change it is an error:

\begin{salamcode}
func main:
    age := 17
    age = 18
    println age
end
\end{salamcode}

\begin{salamout}
error[E013]: cannot reassign immutable variable 'age'; declare it
'mut' to reassign it, or assign to its elements or fields instead
 --> main.salam:3:5
\end{salamout}

\code{E013} in the square brackets is the number of this kind of error;
every kind of error in Salam has its own number, and you can search for it.
The message itself tells you the remedy. If you want the value of a box to
change during the program, write the word \kwd{mut} (short for
``mutable'') before its name when you declare it:

\program{A changeable variable}
\begin{salamcode}
func main:
    mut age := 17
    println "This year:", age
    age = 18
    println "Next year:", age
end
\end{salamcode}

\begin{salamout}
This year: 17
Next year: 18
\end{salamout}

Note the two different signs:
\begin{itemize}
  \item \code{:=} is used only \emph{once}, when the box is made.
  \item \code{=} is for \emph{changing} the value of a box that already
        exists and was declared with \code{mut}.
\end{itemize}
If you write \code{:=} again for a variable you have already made, Salam
reports that the name is already defined.

\begin{note}[Why is Salam so strict?]
You may ask why we should go to the trouble of writing \code{mut}. The
answer is that most values in a program never change: a price, a name, the
number of days in a week. When you declare these without \code{mut}, Salam
guarantees that no part of the program, not even by accident, will change
them. And when you do see \code{mut} somewhere, you know at once: ``this one
changes, keep an eye on it.'' This small habit prevents a great many large
mistakes.
\end{note}

\section{Salam takes every variable seriously}

Another thing you should know right at the start: if you declare a variable
and never use it anywhere, Salam reports an error:

\begin{salamcode}
func main:
    name := "Sara"
    age := 17
    println "Name:", name
end
\end{salamcode}

\begin{salamout}
error[E059]: unused variable 'age': nothing uses it after it is
declared; remove the declaration, or use the value where it was
meant to be used; starting the name with '_' hides this error but
is not advised, since the dead code stays in the program
 --> main.salam:3:5
\end{salamout}

The reason is simple: a variable that is never used is either unnecessary,
or a sign that you forgot something somewhere. In both cases you would
rather know. The remedy is simple too: either use it, or delete its line.

\section{Values are copied}

When you put the value of one variable into another, \emph{a copy} of the
value goes into the second box. The two boxes are not connected:

\program{Copying a value}
\begin{salamcode}
func main:
    mut a := 10
    b := a
    a = 20
    println "a:", a
    println "b:", b
end
\end{salamcode}

\begin{salamout}
a: 20
b: 10
\end{salamout}

Although \code{b} was made from \code{a}, changing \code{a} had no effect on
\code{b}. At the moment of \code{b := a}, the number 10 was copied into the
box \code{b}; from then on each box goes its own way.

\section{Everyday patterns}

There are a few things programmers do with variables every single day. Let
us look at them one by one.

\subsection{The counter}

Perhaps the strangest line for a beginner is this one:

\begin{salamcode}
counter = counter + 1
\end{salamcode}

As a statement of mathematics it is nonsense; no number is equal to one more
than itself. But remember that in programming \code{=} does not mean ``is
equal to''. It means ``put into''. This line says: take the current value of
\code{counter}, add one to it, and put the result back into the same box.
The right-hand side is worked out first, and then the result is stored on
the left.

\program{A counter}
\begin{salamcode}
func main:
    mut counter := 0
    println counter
    counter = counter + 1
    println counter
    counter = counter + 1
    println counter
end
\end{salamcode}

\begin{salamout}
0
1
2
\end{salamout}

This is so common that Salam has shorter ways of writing it.
\code{counter += 1} is exactly \code{counter = counter + 1}, and
\code{counter++} is the same as \code{counter += 1}. All three forms are
correct, and you will see all three in this book:

\begin{salamcode}
func main:
    mut counter := 0
    counter = counter + 1
    counter += 1
    counter++
    println counter
end
\end{salamcode}

\begin{salamout}
3
\end{salamout}

\subsection{Adding up}

The same pattern serves for adding things up. We make a variable with the
value zero and add to it whatever needs to be added:

\program{Adding up a shopping list}
\begin{salamcode}
func main:
    mut total := 0
    total += 4     // bread
    total += 12    // cheese
    total += 3     // milk
    println "Total:", total, "dollars"
end
\end{salamcode}

\begin{salamout}
Total: 19 dollars
\end{salamout}

\subsection{Swapping two values}

Suppose you want to exchange the contents of two boxes. You cannot first put
\code{a} into \code{b} and then \code{b} into \code{a}, because after the
first step the old value of \code{b} is gone. The answer is a third box:

\program{Swapping}
\begin{salamcode}
func main:
    mut a := 5
    mut b := 9
    println "Before the swap:", a, b
    temp := a
    a = b
    b = temp
    println "After the swap:", a, b
end
\end{salamcode}

\begin{salamout}
Before the swap: 5 9
After the swap: 9 5
\end{salamout}

Follow it step by step: \code{temp} keeps a copy of 5, \code{a} receives 9,
and finally \code{b} takes the 5 back from \code{temp}. Notice that we
declared \code{temp} without \code{mut}, because it receives a value once
and never changes.

\section{Constants}

Some values not only stay the same while the program runs, they are fixed by
their very nature: the number of days in a week, a tax rate, the number pi.
For these Salam has the word \kwd{const}:

\program{A constant}
\begin{salamcode}
func main:
    const rate := 9
    price := 500
    tax := (price * rate) / 100
    println "Tax:", tax
    println "Price with tax:", price + tax
end
\end{salamcode}

\begin{salamout}
Tax: 45
Price with tax: 545
\end{salamout}

As far as the program is concerned, \code{const rate := 9} differs little
from \code{rate := 9}; both are unchangeable. The difference is in the
message: \code{const} tells the reader ``this is an important number that
we have put here on purpose.'' Two small points: the name of a constant must
be a single word (unlike variables, which may have several), and constants
may also be declared outside \code{func main}, at the top of the program,
where they are available everywhere. We will do that in Chapter~\ref{ch:functions}.

\begin{deep}[Where is a variable, really?]
The computer's memory really is like a shelf of numbered compartments. Each
compartment has an \textbf{address}, and each holds one small number. When
you write \code{age := 17}, the Salam compiler finds a few free
compartments, ties their address to the name \code{age}, and writes the
number 17 into them. After that, every \code{age} you write is translated
into that address. You never have to see the addresses; names were invented
for exactly this reason. Pieces of text such as \code{"Sara"} take a little
more room (a compartment or so for each letter), but the principle is the
same.
\end{deep}

\section{A more complete program}

Let us write the bill for an online purchase with variables:

\program{A purchase bill}
\begin{salamcode}
func main:
    unit price := 45
    count := 4
    shipping := 30
    items total := unit price * count
    println "Items:", items total
    println "Shipping:", shipping
    println "Amount due:", items total + shipping
end
\end{salamcode}

\begin{salamout}
Items: 180
Shipping: 30
Amount due: 210
\end{salamout}

None of the variables is declared with \code{mut}, because none of them
changes while the program runs. This program could also have been written
with the numbers typed straight in, but then, if the count changed from 4 to
5, we would have had to touch several places and would probably have missed
one.

\begin{summary}
\begin{itemize}
  \item A variable is a named place in memory that holds one value.
  \item \code{name := value} makes a new variable that cannot change
        afterwards. \code{mut name := value} makes one whose value can be
        changed later with \code{=}.
  \item Names may have several words, but they cannot start with a digit
        and cannot be words of the language. Choose descriptive names.
  \item A variable that is declared but never used is an error.
  \item \code{b := a} makes a copy of the value; the two variables are not
        connected.
  \item \code{counter += 1} and \code{counter++} are short forms of
        \code{counter = counter + 1}.
  \item \code{const} is for values that are fixed by nature.
\end{itemize}
\end{summary}

\section*{Exercises}
\markright{Exercises}

\exercise Declare three variables for the length, width and height of a room
(in metres), then work out and print the floor area and the volume of the
room.

\exercise Make a variable called \code{balance} with the starting value
500. Then, in three separate instructions, withdraw 120, deposit 200 and
withdraw 75, printing the balance after each operation. (In a program, write
numbers without thousands separators: \code{1500000}, not
\code{1,500,000}.)

\Needspace{12\baselineskip}
\exercise Trace the program below without running it and say what it
prints. Then run it to check your answer:
\begin{salamcode}
func main:
    mut a := 3
    mut b := a * 2
    a = b + 1
    b = a - b
    println a, b
end
\end{salamcode}

\exercise The program below has two errors. Find and fix them:
\begin{salamcode}
func main:
    score := 10
    score = score + 5
    player name := "Reza"
    println "Score:", score
end
\end{salamcode}

\exercise Make three variables \code{first}, \code{second} and
\code{third} with the values 1, 2 and 3, and without typing any numbers
directly, rotate the values so that \code{first} becomes 2, \code{second}
becomes 3 and \code{third} becomes 1.
